\section{Theoretical GPnn Results}
\label{detailed_proofs}

\label{appx:gpnn}
\subsection{Preliminary results}
    Let \(\rho(\bm{x},\bm{x'})=\sigma_f\sqrt{1-c(\bm{x}/l,\bm{x'}/l)}\) be the kernel-induced distance function over \(\mathbb{R}^d\) (\cite{Scholkopf2001}).
    We define \(\mathbf{x}_{(j,n)}(\bm{x}^*)\) as the \(j^{th}\) nearest neighbour random variable to a test point \(\bm{x}^*\) under \(\rho\), which we abbreviate to \(\mathbf{x}_{(j)}\) when the context is clear, and \(\bm{x}_{(j)}(\bm{x}^*) \in N_m(\bm{x}^*)\) as the realised \(j^{th}\)
    nearest neighbour of the test point \(\bm{x}^*\) from a training set \(X\). From this we define \(\epsilon_i = \rho^2(\mathbf{x}_{(i)},\bm{x}^*)\) and \(\epsilon_{ij} = \rho^2(\mathbf{x}_{(i)},\mathbf{x}_{(j)})\).

\begin{definition}[Support]\label{def:support}
    Let \(P_\mathbf{x}\) be the probability measure of \(\mathbf{x}\) and \(S^\rho_{\bm{x},\epsilon}\) the closed ball of radius \(\epsilon>0\) under the metric \(\rho\) centred at \(\bm{x}\). Then we define \(\support(P_\mathbf{x}) = \{\bm{x} : P_\mathbf{x}(S^\rho_{\bm{x},\epsilon})>0\, \forall\,\epsilon >0\}\).
\end{definition}

\begin{definition}[Weakly-faithful]\label{def:weakly_faithful}
    We define a pair of metrics \(\rho(\cdot,\cdot),\hat{\rho}(\cdot,\cdot)\) to be \emph{weakly-faithful} w.r.t. each other if the following condition holds: The \(m^{th}\) nearest neighbour under \(\hat{\rho}\) converges to the test point as \(n\rightarrow\infty\) if and only if the \(m^{th}\) nearest neighbour under \(\rho\) converges to the test point in the limit.
\end{definition}

\paragraph{Assumptions}
\begin{description}[labelindent=1cm]
\label{desc:assumptions_gpnn}
    \item[(A1)\label{ass:X}] \(\mathbf{x}\overset{\textrm{iid}}{\sim}P_\mathbf{x}\) and \(\mathbf{x}^* \in \support(P_\mathbf{x})\) under the generative metric defined by \(c(\cdot,\cdot)\).
    % \item[(A2)\label{ass:gpnn_k}] \(c(\cdot,\cdot),\hat{c}(\cdot,\cdot)\) are stationary kernels whose induced distance metrics are both \emph{weakly faithful} (\autoref{def:weakly_faithful}) w.r.t. the Euclidean metric. 
    \item[(A2)\label{ass:gpnn_k}] \(c(\cdot,\cdot),\hat{c}(\cdot,\cdot)\) are stationary kernels whose induced distance functions are \emph{weakly faithful} metrics (\autoref{def:weakly_faithful}). 
    \item[(A3)\label{ass:gpnn_y}]  \(y_i=f(\bm{x}_i) + \xi_i\) with \(\xi_i\overset{\textrm{iid}}{\sim}P_\xi\), \(f(\bm{x})\sim \mc{WSRF}(\sigma_f^2 c(./l,./l))\) and \(y_i\given f(\bm{x}_i) \sim P_\xi\) and $\E[\xi]=0$, $\E[\xi^2]=\sigma_\xi^2$.
\end{description}

\textbf{Note:} Assumption \hyperref[{desc:assumptions_gpnn}]{(A2)} is not overly restrictive and encompasses commonly used kernels such as all those mentioned in this paper.

\begin{lemma}
    \label{lem:eps_limit}
    \(\epsilon_i\rightarrow0\) and \(\epsilon_{ij}\rightarrow0\) as \(n\rightarrow\infty\) a.e. with respect to the measure over \(\mathbf{x}\in\mathbb{R}^d\), \(P_\mathbf{x}\), for \(i,j\leq m\), \(\frac{m}{n}\rightarrow0\) and under \hyperref[{desc:assumptions_gpnn}]{(A1-2)}.
\end{lemma}
\begin{proof}
    Lemma 6.1 of \cite{Gyorfi2010} states that \(\norm{\mathbf{x}_{(m,n)}(\bm{x})-\bm{x}}\xrightarrow{n\rightarrow\infty}0\) with probability one (with respect to \(P_\mathbf{x}\)). Their proof can be generalised immediately to state that \(\rho(\mathbf{x}_{(m,n)}(\bm{x}), \bm{x})\xrightarrow{n\rightarrow\infty}0\) by using our definition of support, \ref{def:support}, that directly invokes the metric \(\rho\). Hence \(\epsilon_i\rightarrow 0\) for all \(i\leq m\) (since \(\bm{x}^*\) is in \(\support(P_\mathbf{x})\)). Since \(\rho\) is a metric it satisfies the triangle inequality; hence \(\rho(\mathbf{x}_{(i)},\mathbf{x}_{(j)}) \leq \rho(\mathbf{x}_{(i)},\bm{x}^*) + \rho(\mathbf{x}_{(j)}, \bm{x}^*) \xrightarrow{n\rightarrow\infty}0\) for all \(i,j \leq m\).
\end{proof}

%Finally, by the triangle inequality, \(\norm{\mathbf{x}_{(i)}-\mathbf{x}_{(j)}}\leq \norm{\mathbf{x}_{(i)}-\bm{x}^*} + \norm{\mathbf{x}_{(j)}-\bm{x}^*}\xrightarrow{n\rightarrow\infty}0\) and so \(\epsilon_{ij}\rightarrow0\) for all \(i,j\leq m\).

% \begin{proof}
%     We prove this result in two parts: i) the case of an isotropic kernel and ii) the generalisation to a stationary anisotropic kernel.
%     \begin{enumerate}[label=\roman*)]
%         \item First we prove the convergence of \(\epsilon_i\): Lemma 6.1 of \cite{Gyorfi2010} states that \(\delta_m:=\norm{\mathbf{x}_{(m,n)}(\bm{x})-\bm{x}}\xrightarrow{n\rightarrow\infty}0\) with probability one (with respect to \(P_\mathbf{x}\)). For a stationary isotropic kernel function \(c\), the corresponding metric \(\rho\) can be regarded as a continuous function of \(\delta_m\), \(\rho(\mathbf{x}_{(m)},\bm{x})\) thus (with a slight abuse of notation) \(\lim_{n\rightarrow\infty}\rho^2(\delta_m) = \rho^2(0)\). For such a kernel function, \(c(0)=1\) and so from the definition given above, \(\rho^2(0)=0=\rho(0)\). 
%         Now we can use the result of the first part to show the convergence of \(\epsilon_{ij}\). Since \(\rho\) is a metric it satisfies the triangle inequality; hence \(\rho(\mathbf{x}_{(i)},\mathbf{x}_{(j)}) \leq \rho(\mathbf{x}_{(i)},\bm{x}^*) + \rho(\mathbf{x}_{(j)}, \bm{x}^*) \xrightarrow{n\rightarrow\infty}0\) for all \(i,j \leq m\).
%         \item We now relax this to the anisotropic case by asserting that for a given choice of such a kernel function there exists a constant positive definite matrix \(A\) defining anisotropic distances in the input \(\bm{x}\)-space: \(\delta_m^A := (\mathbf{x}_{(m,n)}(\bm{x})-\bm{x})^TA(\mathbf{x}_{(m,n)}(\bm{x})-\bm{x})\). The metric in this case is a continuous function of \(\delta_m^A\). To complete the proof we use that \(\underline{\lambda(A)}\norm{\mathbf{x}_{(m,n)}(\bm{x})-\bm{x}} \leq (\mathbf{x}_{(m,n)}(\bm{x})-\bm{x})^TA(\mathbf{x}_{(m,n)}(\bm{x})-\bm{x}) \leq \overline{\lambda(A)}\norm{\mathbf{x}_{(m,n)}(\bm{x})-\bm{x}}\), with \(\overline{\lambda},\underline{\lambda}\) the (finite) maximum and minimum eigenvalues respectively, and apply the argument from part i). Note that in the Euclidean case, \(A=I\), the identity.
%     \end{enumerate}
% \end{proof}

\begin{lemma}
    \label{lem:gpnn_limit}
    For an \(m\)-GPnn under the assumptions \hyperref[{desc:assumptions_gpnn}]{(A1-3)},
    \[\lim_{n\rightarrow\infty} \kstar_N^TK_N^{-1}\kstar_N = \sigma_f^2-\sigma_\xi^2m^{-1} + \bigO{m^{-2}}.\]
\end{lemma}
\begin{proof}
    From \autoref{lem:eps_limit} we have that 
    \(\lim_{n\rightarrow\infty}k(\mathbf{x}_{(j)}(\bm{x}^*),\bm{x}^*) =
    \lim_{n\rightarrow\infty} (\sigma_f^2 - \epsilon_i) = \sigma_f^2\) and \(\lim_{n\rightarrow\infty} k(\mathbf{x}_{(i)}(\bm{x}^*), \mathbf{x}_{(j)}(\bm{x}^*)) = \lim_{n\rightarrow\infty}(\sigma_f^2 - \epsilon_{ij}) = \sigma_f^2\). 
    As a result, \(\kstar_N\rightarrow \sigma_f^2 \bm{1}\) and 
\begin{align}
    K^\infty:= \lim_{n\rightarrow\infty}K_N & = \sigma_\xi^2 I + \sigma_f^2\bm{11}^T.
\end{align}
Now using Sherman-Morrison and the continuity of matrix inverse and matrix-matrix products:
\begin{align}
    (A + \bm{bc}^T)^{-1} &= A^{-1} - \frac{A^{-1}\bm{bc}^TA^{-1}}{1+\bm{c}^TA^{-1}\bm{b}} \label{eq:sherman-morrison}\\
    (K^\infty)^{-1} = (\sigma_\xi^2I + \sigma_f^2 \bm{11}^T)^{-1} &= \frac{1}{\sigma_\xi^2}\left(I -
    \sigma_f^2\frac{\bm{11}^T}{\sigma_\xi^2 + \sigma_f^2\bm{1}^T\bm{1}}\right) \label{eq:lim_Kinv}\\
    \bm{1}^T(K^\infty)^{-1}\bm{1} &= \frac{m}{\sigma_\xi^2}\left(1 -
    \frac{m\sigma_f^2}{\sigma_\xi^2 + m\sigma_f^2}\right) \nonumber\\
    &= \frac{m}{\sigma_\xi^2}\left(1-m\sigma_f^2\frac{1}{m\sigma_f^2}\left(1-\frac{\sigma_\xi^2}{m\sigma_f^2}
    +\frac{\sigma_\xi^4}{m^2\sigma_f^4}-\bigO{m^{-3}}\right)\right) \nonumber\\
    &= \frac{1}{\sigma_f^2} -\frac{\sigma_\xi^2}{m\sigma_f^4} + \bigO{m^{-2}} \label{eq:Kinf_sum}. 
\end{align}
Thus,
\begin{align}
    \lim_{n\rightarrow\infty}\kstar_N^TK_N^{-1}\kstar_N &=
    \sigma_f^4\bm{1}^T(K^\infty)^{-1}\bm{1} = \sigma_f^2-\sigma_\xi^2m^{-1} + \bigO{m^{-2}}.
\end{align}
\end{proof}

\begin{lemma}[\(\mc{WSRF}\) expectations]\label{lem:wsrf_exp}
    Under \hyperref[ass:gpnn_y]{(A3)}, \(\E_{\mathbf{y},\mathrm{y^*}}\{\bm{y}y^*\} = \kstar\) and \(\E_{\mathbf{y}}\{\bm{yy}^T\} = K\).
\end{lemma}
\begin{proof}
    By assumption on the covariance properties of \(y\) and the independence and zero-mean of the additive noise, \(\E_\mathbf{y}\{y_iy_j\}=k(\bm{x}_i,\bm{x}_j)\). Extending this to the joint distribution over \(\bm{y},y^*\) is straightforward and gives the results stated.
\end{proof}
\autoref{lem:wsrf_exp} is subsequently assumed to be in use throughout \ref{sec:limit_proofs}.

\subsection{Limit proofs}
\label{sec:limit_proofs}
In the following statements only misspecification of type (d) and/or (e) (\autoref{misspec}) is considered to be at work.

% Additionally, we note that for functions of the test observation only (\(f(y^*)\)), when we take expectations over \(\mathbf{y}\), they are equivalent to expectations over \(\mathbf{y}_N\) so that \(\Ey{f(y^*)}=\E_\mathbf{y}\{\E_{y^*\given\mathbf{y}}\{f(y^*)\}\}=\E_{\mathbf{y}_N}\{\E_{y^*\given\mathbf{y}_N}\{f(y^*)\}\}\). We justify this claim by the fact that the function \(y(\cdot)\) is distributed according to a stochastic process, by assumption, and hence by Kolmogorov's consistency theorem (see e.g. \cite{athreya2006measure}) this implies that any finite set of marginal observations satisfies the corresponding consistency conditions such that the above is true.

\begin{lemma}[MSE limit]
    \label{lem:mse_limit}
    Under the assumptions \hyperref[{desc:assumptions_gpnn}]{(A1-3)}, for fixed \(m<\infty\), the predictive GPnn given in \autoref{exact_NN_prediction} converges pointwise in the sense of MSE wrt \(P_\mathbf{x}\)-a.e. as
    \[\lim_{n\rightarrow\infty}f_n^{\MSE}(\bm{\theta}) = \sigma_\xi^2(1+m^{-1}) - \bigO{m^{-2}}.\]
\end{lemma}
\begin{proof}\label{proof:mse_limit}
    This follows from \autoref{lem:gpnn_limit} by expanding the definition
    of \(\MSE\):
    \begin{align*}
        \lim_{n\rightarrow\infty}f_n^{\MSE}(\bm{\theta})
        &=\lim_{n\rightarrow\infty}\Ey{\abs{y^* - \mu_N^*}^2} \\
        % &= \lim_{n\rightarrow\infty}\E_\mathbf{y}\E_\mathrm{y^*}\left[{y^*}^2 +
        % (\mu_N^*)^2 - 2y^*\mu_N^*\given \mathbf{y}\right]\\
        % &= \lim_{n\rightarrow\infty}\left[\E_\mathrm{y^*}\{{y^*}^2\} + \E_\mathbf{y}\{{\mu^*_N}^2\} - 2\E_\mathbf{y}\{\E_{\mathrm{y^*}\given\mathbf{y}}\{y^*\mu^*_N\}\}\right]\\
        &= \lim_{n\rightarrow\infty}\left[\E_\mathrm{y^*}\{{y^*}^2\} + \E_\mathbf{y}\{{\mu^*_N}^2\} - 2\E_{\mathbf{y},\mathrm{y^*}}\{\kstar_N^TK_N^{-1}\bm{y}_Ny^*\}\right]\\
        &= \sigma_f^2 + \sigma_\xi^2 -
        \lim_{n\rightarrow\infty}\E_\mathbf{y}\{{\mu^*_N}^2\} \\
        &= \sigma_\xi^2(1 +m^{-1}) - \bigO{m^{-2}}.
    \end{align*}
    Since \(\E_\mathbf{y}\{{\mu^*_N}^2\}=\E_\mathbf{y}\{\kstar_N^TK_N^{-1}\bm{y}_N\bm{y}_N^TK_N^{-1}\kstar_N\}=\kstar_N^TK_N^{-1}\kstar_N\), and by assumption \(\E_{\mathbf{y},\mathrm{y^*}}\{\bm{y}_Ny^*\}=\kstar_N\), even under a \(\mc{WSRF}\) generative model (\autoref{lem:wsrf_exp}).
\end{proof}

\begin{corollary}[NLL limit]
\label{cor:nll_limit}
    \[
        \lim_{n\rightarrow\infty}f_n^{\NLL}(\bm{\theta}) = \half\log(\sigma_\xi^2(1+m^{-1})) +\half + \half\log2\pi - \bigO{m^{-2}}.
    \]
\end{corollary}
\begin{proof}\label{proof:nll_limit}
    The proof follows straightforwardly from \autoref{lem:gpnn_limit} and because \(\predvar_N = \sigma_f^2 + \sigma_\xi^2 - \kstar_N^TK_N^{-1}\kstar_N\).
    \begin{align*}
        2\Ey{l^*_N} &= \Ey{\log\predvar_N + \frac{(y^*-\mu_N^*)^2}{\predvar_N} + \log2\pi } \\
        &= \log\predvar_N + 1 + \log2\pi \\
        \lim_{n\rightarrow\infty}2\Ey{l^*_N} &= \log(\sigma_f^2 + \sigma_\xi^2 -(\sigma_f^2 - \sigma_\xi^2 m^{-1} + \bigO{m^{-2}})) + 1 + \log2\pi \\
        &= \log\left(\sigma_\xi^2(1+m^{-1}) - \bigO{m^{-2}}\right) +1 + \log2\pi \\
        &= \log\sigma_\xi^2 + m^{-1} + 1 + \log2\pi - \bigO{m^{-2}}.
    \end{align*}
\end{proof}

\subsubsection{Full misspecification}
For the remainder of \ref{sec:limit_proofs} we assume that the full range of possible misspecifications ((a)-(e)) outlined in \autoref{misspec} are in action. We refer to this case as ``fully-misspecified'' and introduce the notation \(\hat{\mu}_N^*,\predvar[\hat]_N\) to be understood to mean the predictive mean and variance under these misspecifications.
\begin{lemma}[Fully misspecified MSE limit]
    \label{lem:mismatched_mse_limit}
    For a fully misspecified model, asymptotically
    \[\lim_{n\rightarrow\infty}f_n^{\MSE}(\hat{\bm{\theta}}) = \sigma_\xi^2(1+m^{-1}) \pm \bigO{m^{-2}}.\]
    provided the misspecified kernel distance metric is \emph{weakly faithful} in the sense that the \(m^{th}\) nearest neighbour converges under both the true and misspecified metrics (\autoref{def:weakly_faithful}).
\end{lemma}
\begin{proof}\label{proof:mismatched_mse_limit}
\begin{align*}
        \E_\mathbf{y}\bigg\{\E_\mathrm{y^*}\left[(y^* - \hat{\mu}_N^*)^2 \given \mathbf{y}\right] \bigg\}
        &= \E_\mathbf{y}\bigg\{\E_\mathrm{y^*}\left[{y^*}^2 - 2y^*\hat{\mu}_N^* + (\hat{\mu}_N^*)^2\given \mathbf{y}\right]\bigg\} \\
        &= \E_\mathbf{y}\left\{\predvar_N + \mu^{*\,2}_N - 2\mu^*_N\hat{\mu}_N^* + (\hat{\mu}_N^*)^2 \right\} \\
        &= \underbrace{\predvar_N}_{(a)} + \underbrace{\kstar_N^TK_N^{-1}{\kstar_N}}_{(b)} - 2\underbrace{\kstar_N^T\hat{K}_N^{-1}\kstar[\hat]_N}_{(c)} + \underbrace{\kstar[\hat]_N^T\hat{K}_N^{-1}K_N\hat{K}_N^{-1}\kstar[\hat]_N}_{(d)}.
    \end{align*}
    We can use standard results to state that \((a)+(b)=\sigma_f^2+\sigma_\xi^2\).
    Then we define \(\hat{\gamma}=\frac{\hat{\sigma}_f^2}{\hat{\sigma}_\xi^2+m\hat{\sigma}_f^2}\) and expand it in terms of \(m^{-1}\):
    \[1-m\hat{\gamma} = \frac{\hat{\sigma}_\xi^2}{m\hat{\sigma}_f^2} - \frac{\hat{\sigma}_\xi^4}{m^2\hat{\sigma}_f^4} + \bigO{m^{-3}}. \] 
    In a manner similar to \autoref{lem:gpnn_limit} we use this result to compute:
    \begin{align*}
        \lim_{n\rightarrow\infty}(c) &= \sigma_f^2\bm{1}^T\hat{\sigma}_\xi^{-2}(I-\hat{\gamma}\bm{11}^T)\bm{1}\hat{\sigma}_f^2 \\
        &=\frac{\sigma_f^2\hat{\sigma}_f^2}{\hat{\sigma}_\xi^2}m(1-m\hat{\gamma})\\
        &=\frac{\sigma_f^2\hat{\sigma}_f^2}{\hat{\sigma}_\xi^2}\left(\frac{\hat{\sigma}_\xi^2}{\hat{\sigma}_f^2} - \frac{\hat{\sigma}_\xi^4}{m\hat{\sigma}_f^4}\right) + \bigO{m^{-2}} \\
        &= \sigma_f^2 - \frac{\sigma_f^2\hat{\sigma}_\xi^2}{m\hat{\sigma}_f^2} + \bigO{m^{-2}}
    \end{align*}
    and
    \begin{align*}
        \lim_{n\rightarrow\infty}(d) &= \frac{\hat{\sigma}_f^4}{\hat{\sigma}_\xi^4}\bm{1}^T(I-\hat{\gamma}\bm{11}^T)(\sigma_\xi^2I+\sigma_f^2\bm{11}^T)(I-\hat{\gamma}\bm{11}^T)\bm{1}\\
        &=\frac{\hat{\sigma}_f^4}{\hat{\sigma}_\xi^4}\bm{1}^T\left[\sigma_\xi^2I + \sigma_f^2\bm{11}^T - 2\sigma_\xi^2\hat{\gamma}\bm{11}^T + \hat{\gamma}^2\sigma_\xi^2m\bm{11}^T -2\sigma_f^2\hat{\gamma}m\bm{11}^T + \sigma_f^2\hat{\gamma}^2m^2\bm{11}^T\right]\bm{1}\\
        &= \frac{\hat{\sigma}_f^4}{\hat{\sigma}_\xi^4}m(\sigma_\xi^2+m\sigma_f^2)\left[1-2m\hat{\gamma}+m^2\hat{\gamma}^2\right] \\
        &= \frac{\hat{\sigma}_f^4}{\hat{\sigma}_\xi^4}m(\sigma_\xi^2+m\sigma_f^2)(1-m\hat{\gamma})^2 \\
        &= \frac{\hat{\sigma}_f^4}{\hat{\sigma}_\xi^4}m(\sigma_\xi^2+m\sigma_f^2)\left(\frac{\hat{\sigma}_\xi^4}{m^2\hat{\sigma}_f^4} - 2\frac{\hat{\sigma}_\xi^6}{m^3\hat{\sigma}_f^6} + \bigO{m^{-4}}\right) \\
        &= \sigma_f^2 + \frac{\sigma_\xi^2}{m} - 2\frac{\sigma_f^2}{\hat{\sigma}_f^2}\frac{\hat{\sigma}_\xi^2}{m} \pm \bigO{m^{-2}},
    \end{align*}
    where we have used the expansion of \(1-m\hat{\gamma}\) given earlier. Putting these results together gives
    \begin{align*}
        \lim_{n\rightarrow\infty}f_n^{\MSE}(\hat{\bm{\theta}}) &= \lim_{n\rightarrow\infty} [(a) + (b) -2 (c) + (d)] \\
        &= \sigma_f^2 + \sigma_\xi^2 - 2\left(\sigma_f^2 - \frac{\sigma_f^2\hat{\sigma}_\xi^2}{m\hat{\sigma}_f^2}\right) + \sigma_f^2 + \frac{\sigma_\xi^2}{m} - 2\frac{\sigma_f^2}{\hat{\sigma}_f^2}\frac{\hat{\sigma}_\xi^2}{m} \pm \bigO{m^{-2}} \\
        &= \sigma_\xi^2(1+m^{-1}) \pm \bigO{m^{-2}}.
    \end{align*}
\end{proof}

\begin{lemma}[Calibration limit under full misspecification]
    \label{lem:mismatched_cal_limit}
    \begin{align*}
    \lim_{n\rightarrow\infty} f_n^{\CAL}(\hat{\bm{\theta}}) &= \frac{\sigma_\xi^2}{\hat{\sigma}_\xi^2} \pm \bigO{m^{-2}}.
    \end{align*}
\end{lemma}
\begin{proof}\label{proof:mismatched_cal_limit}
        We use continuity to write
    \begin{align*}
        \lim_{n\rightarrow\infty}\Ey{\frac{(y^*-\hat{\mu}_N^*)^2}{\predvar[\hat]_N}}  &= \left(\lim_{n\rightarrow\infty}\frac{1}{\predvar[\hat]_N}\right)\left(\lim_{n\rightarrow\infty}f_n^{\MSE}(\hat{\bm{\theta}})\right).
    \end{align*}
    By direct application of \autoref{lem:gpnn_limit} \(\predvar[\hat]_N\xrightarrow{n\rightarrow\infty} \hat{\sigma}_\xi^2(1+m^{-1}) - \bigO{m^{-2}}\) and thus
    \[\lim_{n\rightarrow\infty} f_n^{\CAL}(\hat{\bm{\theta}}) = \frac{\sigma_\xi^2}{\hat{\sigma}_\xi^2} \pm \bigO{m^{-2}}.\]
\end{proof}

\begin{corollary}[NLL limit under full misspecification]
    \label{cor:mismatched_nll_limit}
    \begin{align*}
    \lim_{n\rightarrow\infty}f_n^{\NLL}(\hat{\bm{\theta}}) &= \half\log\left(\hat{\sigma}_\xi^2(1 + m^{-1})\right) + \half\frac{\sigma_\xi^2}{\hat{\sigma}_\xi^2} +\half\log2\pi \pm \bigO{m^{-2}}. 
    \end{align*}
\end{corollary}
\begin{proof}\label{proof:mismatched_nll_limit}
We start with
    \begin{align*}
        2f_n^{\NLL}(\hat{\bm{\theta}}) &= \Ey{\log\predvar[\hat]_N + \frac{(y^*-\hat{\mu}_N^*)^2}{\predvar[\hat]_N} + \log2\pi}.
    \end{align*}
        For the second term we use \autoref{lem:mismatched_cal_limit} so that we have 
    \begin{align*}
        \lim_{n\rightarrow\infty}2f_n^{\NLL}(\hat{\bm{\theta}}) &= \log\hat{\sigma}_\xi^2 +  m^{-1} + \frac{\sigma_\xi^2}{\hat{\sigma}_\xi^2} +\log2\pi \pm \bigO{m^{-2}}.
    \end{align*}
\end{proof}

\begin{proof}[Proof of \autoref{thm:gpnn_limit}]\label{proof:gpnn_limit}
    We construct the proof using all of the intermediate results given above. In particular \autoref{eq:mse_lim} follows from \autoref{lem:mismatched_mse_limit}, \autoref{eq:cal_lim} from \autoref{lem:mismatched_cal_limit} and \autoref{eq:nll_lim} from \autoref{cor:mismatched_nll_limit}.
\end{proof}

